\chapter{Fundamentação Teórica}\label{chp:fund}

Esta seção apresenta uma breve fundamentação téorica sobre.......

\section{Segurança e Vazamento de Dados no Desenvolvimento Assistido por IA}

A adoção de ferramentas de desenvolvimento assistidas por modelos de linguagem de grande escala introduz uma nova superfície de exposição de dados: cada interação entre o desenvolvedor e o assistente, seja um trecho de código, uma mensagem de commit ou o contexto de um arquivo aberto no editor é enviada a um provedor externo de inteligência artificial para ser processada. Esse fluxo caracteriza o que a literatura denomina \textit{data leakage}, definido como a transmissão não autorizada ou não intencional de informação sensível para fora do domínio de controle da organização, e sua variação mais específica ao contexto deste trabalho, o \textit{code leakage}, que se refere à exposição de código-fonte proprietário, segredos de aplicação e credenciais de acesso durante o uso de ferramentas externas de desenvolvimento.

Diferentemente de vazamentos causados por falhas de configuração de infraestrutura ou ataques direcionados, a exposição de dados em ferramentas assistidas por LLM ocorre como efeito colateral do uso legítimo da ferramenta: o dado sensível não é roubado, é enviado voluntariamente pelo próprio usuário como parte do fluxo normal de trabalho, o que torna esse risco mais difícil de mitigar por controles de segurança perimetrais tradicionais, uma vez que o tráfego se origina de dentro do ambiente de desenvolvimento e é direcionado a um destino legítimo e autorizado pela organização.

Como o ocorrido na empresa \textit{Samsung} em 2023, onde os seus engenheiros usaram o \textit{ChatGPT} para avaliar o código-fonte e otimizar as sequências de testes para identificar falhas nos chips que estavam projetando. Causando três situações de vazamento de dados no intervalo de um mês, culminando na proibição do uso da ferramenta pelos funcionários da empresa e na discussão de formas seguras de utilização.

Em 2026, a ferramenta de detecção de segredos \textit{GitGuardian}, da plataforma \textit{GitHub}. Registrou aproximadamente 29 milhões de chaves, tokens e credenciais expostas. E mais de 24.000 segredos únicos em arquivos de configuração MCP no GitHub público. Reforçando o fato de que o uso irrestrito das ferramentas trazem riscos a segurança. 


\section{Modelos de Linguagem de Grande Escala (LLMs) e Riscos de Privacidade}

Modelos de linguagem de grande escala (LLMs) são sistemas de aprendizado profundo treinados sobre grandes volumes de texto, capazes de gerar e interpretar linguagem natural e código-fonte a partir de um contexto (\textit{prompt}) fornecido pelo usuário. No cenário de ferramentas de desenvolvimento assistidas por IA, esse contexto tipicamente inclui não apenas a instrução explícita do usuário, mas também trechos do código-fonte do projeto, estrutura de arquivos e, em alguns casos, o histórico de interações anteriores, dados que trafegam até a infraestrutura do provedor do modelo, frequentemente hospedada por terceiros e fora do domínio de controle direto da organização que utiliza a ferramenta.

\citeonline{xu2025privacypreserving} sistematizam os principais vetores de risco à privacidade decorrentes desse fluxo, categorizando-os em riscos durante o treinamento do modelo, quando dados sensíveis podem ser memorizados e posteriormente reproduzidos, e riscos durante a inferência, quando o próprio conteúdo do prompt enviado pelo usuário constitui a superfície de exposição. Os autores destacam que técnicas de computação segura, como criptografia homomórfica, ambientes de execução confiável (TEE) e computação multipartidária segura (MPC), têm sido propostas para mitigar esses riscos, ainda que com custos computacionais e de latência relevantes para cenários de uso interativo. Em linha semelhante, \citeonline{zhang2024privacyasst} argumentam que, em arquiteturas de agentes que utilizam ferramentas externas, o risco de exposição se amplia porque o agente frequentemente repassa dados do usuário não apenas ao modelo de linguagem, mas também às ferramentas que ele orquestra como motores de busca, interpretadores de código e integrações com sistemas de terceiros, multiplicando os pontos de possível vazamento.

\section{Assistentes de Desenvolvimento (Code Assistants)}

Assistentes de desenvolvimento baseados em IA (\textit{code assistants}) são ferramentas integradas ao ambiente de desenvolvimento (IDE) ou à linha de comando que utilizam LLMs para auxiliar tarefas como completude de código, geração de funções, refatoração e revisão. Para operar de forma útil, esses assistentes precisam de acesso ao contexto do projeto, arquivos abertos, estrutura de diretórios, histórico de versionamento, e esse contexto é, em arquiteturas convencionais baseadas em provedores de nuvem, transmitido integralmente ao provedor externo a cada interação.

\citeonline{elatrassi2025codewisp} ilustram uma abordagem arquitetural alternativa a esse modelo com o \textit{CodeWisp}, um assistente de código que opera localmente, utilizando geração aumentada por recuperação (RAG) com indexação semântica guiada por árvore de sintaxe abstrata (AST) para localizar e fornecer apenas os trechos de código relevantes ao contexto da tarefa, sem depender de um LLM hospedado externamente. Embora a proposta demonstre viabilidade técnica, com recall de 0,96 e \textit{Mean Reciprocal Rank} de 0,95 em tarefas de completude, ela pressupõe a substituição do provedor externo por um modelo local, o que não se aplica a cenários em que o uso de LLMs de terceiros (como os grandes provedores comerciais) já está consolidado como decisão de produto ou de negócio da organização.

\section{Detecção de Dados Sensíveis}

A detecção de dados sensíveis é o processo de identificar, em um fluxo de texto ou código, elementos que representam informação confidencial, credenciais de acesso, chaves de API, dados pessoalmente identificáveis (PII) ou segredos de aplicação, antes que esse conteúdo seja transmitido a um destino não confiável. As abordagens de detecção evoluíram de técnicas baseadas puramente em expressões regulares e análise de entropia, que apresentam alta taxa de falsos positivos por não considerarem o contexto semântico do conteúdo analisado, para abordagens híbridas e baseadas em aprendizado de máquina.

\citeonline{rahman2025secret} demonstram essa evolução ao combinar extração de candidatos via regex com classificação por modelos de linguagem ajustados (LLaMA-3.1 8B e Mistral-7B), alcançando F1-score de 0,9852 na tarefa de classificação binária de segredos em código-fonte. \citeonline{tsebak2025enhancing} seguem abordagem semelhante ao integrar um modelo transformer contextual (CodeBERT) a um classificador de gradient boosting (XGBoost), reportando redução significativa de falsos positivos frente a métodos tradicionais de entropia. No contexto de dados pessoais em vez de segredos técnicos, \citeonline{sakapertana2025advancing} comparam algoritmos clássicos de aprendizado supervisionado, SVM e Naive Bayes para detecção de PII em respostas de API REST, com SVM atingindo 80,96\% de acurácia, enquanto \citeonline{khan2025aipowered} propõem o uso de modelos generativos em regime \textit{zero-shot}, capazes de se adaptar a novos critérios de sensibilidade sem necessidade de reentreinamento, uma característica relevante para cenários em que a definição do que constitui "dado sensível" pode variar entre organizações.

\section{Mitigação: Sanitização e Redação de Dados}

Uma vez detectado, o dado sensível precisa ser tratado antes de prosseguir ao destino externo. A literatura de privacidade descreve um espectro de técnicas de mitigação: a \textit{redação} (\textit{redaction}), que consiste na remoção ou ocultação total do dado sensível identificado; a \textit{sanitização} (\textit{sanitization}), termo mais amplo que engloba qualquer transformação do conteúdo para reduzir seu risco antes da transmissão; e a \textit{anonimização} e \textit{pseudonimização}, que substituem o valor original por um identificador não reversível ou reversível apenas sob controle específico, respectivamente, permitindo preservar a utilidade estrutural do dado sem expor seu valor real.

\citeonline{zhang2024privacyasst} aplicam esse princípio no nível arquitetural, propondo um mecanismo de embaralhamento (\textit{shuffling}) de atributos combinado a criptografia homomórfica para que agentes de LLM operem sobre dados protegidos sem acesso ao valor original do usuário. \citeonline{pan2025selective} tratam a mitigação em outro ponto do ciclo de vida do modelo, aplicando ruído customizado a tokens privados durante o processo de ajuste fino (\textit{fine-tuning}) em cenários de aprendizado federado, de forma seletiva, isto é, calibrando o nível de proteção conforme a sensibilidade de cada token, em vez de aplicar uma transformação uniforme a todo o conteúdo. Essa noção de seletividade é particularmente relevante para o mecanismo proposto neste trabalho, pois a redação indiscriminada de todo o tráfego comprometeria a utilidade do assistente de código para o desenvolvedor.

\section{Arquiteturas Intermediárias de Proteção}

Arquiteturas de proteção posicionadas entre o cliente e o serviço externo, como \textit{proxies} reversos, \textit{API gateways} e \textit{middlewares}, constituem um padrão consolidado em segurança de aplicações para implementar pontos de aplicação de política (\textit{policy enforcement points}), nos quais regras de segurança, autenticação e inspeção de conteúdo podem ser aplicadas de forma centralizada e transparente às aplicações que consomem o serviço, sem exigir alterações no código-cliente.

\section{Avaliação de Impacto Operacional}

A introdução de qualquer mecanismo de inspeção ou transformação de tráfego entre o cliente e o provedor externo tem potencial de impactar a experiência de uso da ferramenta, seja por meio do aumento de latência introduzido pelo processamento adicional, seja pela redução de utilidade quando o mecanismo de proteção remove ou distorce informação relevante ao funcionamento correto do assistente. A avaliação desse impacto operacional, portanto, não pode se limitar a métricas de eficácia da detecção como precisão e revocação, devendo incorporar também métricas de eficiência, como tempo de resposta e taxa de falsos positivos, que representam o custo prático de se utilizar o mecanismo de proteção no dia a dia do desenvolvedor.

Os trabalhos revisados na Seção 3.4 fornecem parâmetros de referência para essa avaliação: \citeonline{rahman2025secret} e \citeonline{tsebak2025enhancing} reportam F1-score e taxa de falsos positivos como métricas centrais de eficácia de seus classificadores, enquanto a redução de falsos positivos é tratada por \citeonline{tsebak2025enhancing} como requisito prático para viabilizar a adoção da ferramenta em fluxos de trabalho reais, um paralelo direto ao critério de usabilidade que este trabalho pretende avaliar para o mecanismo de proxy proposto. Nenhum dos trabalhos revisados, no entanto, avalia explicitamente a métrica de latência introduzida pelo mecanismo de proteção, o que reforça a avaliação de impacto operacional como uma contribuição relevante deste trabalho frente à literatura atual.

\section{Governança de Dados e IA}

Governança de dados refere-se ao conjunto de políticas, papéis, processos e controles que definem como os dados de uma organização são geridos, protegidos e utilizados ao longo de seu ciclo de vida. No contexto de adoção de ferramentas de IA generativa no desenvolvimento de software, a governança se materializa na definição de quais categorias de dados podem ou não ser enviadas a provedores externos, em quais condições, e sob qual mecanismo de controle e auditoria essa política é efetivamente aplicada, e não apenas declarada.

\citeonline{xu2025privacypreserving} situam esse desafio ao observar que técnicas de computação segura, embora tecnicamente eficazes, ainda carecem de mecanismos maduros de integração com políticas de governança organizacionais em escala, o que limita sua adoção prática fora de ambientes controlados de pesquisa. \citeonline{pan2025selective} ilustram, por outro lado, como decisões de governança podem ser incorporadas diretamente ao mecanismo técnico de proteção, ao permitir que o nível de proteção aplicado a cada dado seja calibrado conforme sua sensibilidade, uma abordagem que aproxima a implementação técnica da política de governança que a motiva, princípio que orienta o desenho do mecanismo de detecção e redação proposto neste trabalho.
    
\section{Trabalhos Relacionados}
\label{sec:trabalhos_relacionados}

Esta seção apresenta uma análise comparativa dos trabalhos mais próximos ao tema desta pesquisa, identificados a partir da revisão sistemática conduzida com as Strings A a E (\ref{sec:strings_busca}). Os trabalhos foram organizados em três frentes: detecção de dados sensíveis, técnicas de preservação de privacidade em modelos de linguagem, e arquiteturas alternativas para redução de exposição de dados.

\subsection{Detecção de Dados Sensíveis}

Um conjunto expressivo dos trabalhos analisados concentra-se na detecção de segredos, credenciais e informações pessoalmente identificáveis (PII) em artefatos de software. \citeonline{rahman2025secret} propõem uma abordagem híbrida que combina extração de candidatos via expressões regulares com classificação por meio de modelos de linguagem ajustados (LLaMA-3.1 8B e Mistral-7B), alcançando F1-score de 0,9852 na classificação binária e acurácia de 98,2\% na classificação multiclasse, superando abordagens baseadas exclusivamente em regex. Em linha semelhante, \citeonline{tsebak2025enhancing} combinam um modelo transformer contextual (CodeBERT) com um classificador de gradient boosting (XGBoost) para detecção de segredos em sistemas de controle de versão, reduzindo significativamente a taxa de falsos positivos frente a métodos tradicionais baseados em entropia.

No contexto de APIs REST, \citeonline{sakapertana2025advancing} comparam algoritmos clássicos de aprendizado de máquina, Support Vector Machines e Naive Bayes para detecção de PII em respostas de API, obtendo acurácia de 80,96\% com SVM. Já \citeonline{khan2025aipowered} exploram modelos generativos em regime \textit{zero-shot} para classificação de conteúdo sensível, destacando a vantagem de adaptação dinâmica a novos critérios de sensibilidade sem a necessidade de reentreinamento.

Observa-se que, embora estes trabalhos apresentem resultados robustos em precisão de detecção, todos são avaliados sobre artefatos estáticos, código-fonte versionado, respostas de API previamente coletadas ou documentos e não sobre tráfego de rede interceptado em tempo real entre uma ferramenta de desenvolvimento e um provedor externo de inteligência artificial.

\subsection{Preservação de Privacidade em Modelos de Linguagem}

Uma segunda frente de pesquisa trata da preservação de privacidade no próprio ciclo de vida do modelo de linguagem. \citeonline{zhang2024privacyasst} propõem o \textit{PrivacyAsst}, um framework voltado a agentes de LLM que utilizam ferramentas externas, empregando criptografia homomórfica e um mecanismo de embaralhamento de atributos para ocultar entradas individuais dos usuários. \citeonline{pan2025selective} tratam do problema em outro ponto do ciclo de vida do modelo, propondo um mecanismo de aprendizado federado seletivo que adiciona ruído customizado a tokens privados durante o ajuste fino (\textit{fine-tuning}) de LLMs, mitigando a exposição de dados de domínio específico. \citeonline{xu2025privacypreserving} sistematizam, em um trabalho de revisão, as principais técnicas de computação segura aplicadas a LLMs, criptografia homomórfica, ambientes de execução confiável (TEE) e computação multipartidária segura (MPC), consolidando o estado da arte e apontando a escalabilidade como um desafio ainda em aberto.

Esses trabalhos convergem para soluções aplicadas ao \textit{treinamento} ou à \textit{arquitetura interna do agente}, e não a uma camada de interceptação independente, capaz de operar de forma agnóstica em relação à ferramenta de desenvolvimento e ao provedor de IA utilizados.

\subsection{Arquiteturas Alternativas de Redução de Exposição}

Uma abordagem distinta é adotada por \citeonline{elatrassi2025codewisp}, que apresentam o \textit{CodeWisp}, um assistente de código local baseado em geração aumentada por recuperação (RAG), com indexação semântica guiada por árvore de sintaxe abstrata (AST). O sistema alcança recall de 0,96 e \textit{Mean Reciprocal Rank} de 0,95 em tarefas de completude de código, evitando a necessidade de envio de dados a provedores de LLM hospedados em nuvem. Trata-se, no entanto, de uma estratégia de \textit{eliminação} do provedor externo, e não de \textit{sanitização} do tráfego já direcionado a ele, uma premissa que não se aplica a cenários em que o uso de provedores externos de IA é uma decisão de negócio já estabelecida e não passível de substituição imediata por um modelo local.

\subsection{Síntese e Posicionamento do Trabalho}

A análise dos oito trabalhos revisados evidencia uma lacuna: nenhum deles combina, simultaneamente, (i) a interceptação de tráfego HTTP em tempo real entre ferramentas de desenvolvimento assistidas por LLM e provedores externos de IA, (ii) o agnosticismo em relação à ferramenta cliente e ao provedor de destino, e (iii) a avaliação sistemática do impacto operacional,latência e utilidade da camada de proteção introduzida. Os trabalhos voltados à detecção (\citeonline{rahman2025secret}; \citeonline{sakapertana2025advancing}; \citeonline{khan2025aipowered}; \citeonline{tsebak2025enhancing}) concentram-se em artefatos estáticos; os voltados à preservação de privacidade em LLM (\citeonline{zhang2024privacyasst}; \citeonline{pan2025selective}; \citeonline{xu2025privacypreserving}) atuam no treinamento ou na arquitetura do agente; e a alternativa arquitetural identificada (\citeonline{elatrassi2025codewisp}) resolve o problema por substituição do provedor, não por interceptação. É nesse espaço que se posiciona o mecanismo de proxy reverso proposto neste trabalho, ao aplicar técnicas de detecção de dados sensíveis diretamente sobre o tráfego HTTP em trânsito, de forma transparente à ferramenta de desenvolvimento e ao provedor de IA utilizados.




